\documentclass{atomic}

\title{Boundaries}
\newcommand{\breadcrumb}{Topology $>$ Topological Operations}

\begin{document}
    \pagestyle{plain} \currentdoc{note} \noindent
    The \textit{boundary} $\partial S$ of $S \subseteq X$ is its \ex{Closures}{closure} without its \ex{Interiors}{interior}:
    \begin{align*}
        \partial S &= \left( \inte S \cup \ext S \right)^c \\
        &= \overline{S} \setminus \inte S. && (1)
    \end{align*}

    \begin{itemize}
        \item $\left( \partial S \right)^c \in \mathcal{T}$
    \end{itemize}
    
    \begin{proof}[1]
        \phantom{}\vspace{-2em}
        \begin{align*}
            \partial S &= \left( \text{int}\,S \cup \text{ext}\, S \right)^c \\
                       &= \left( \text{int}\,S \right)^c \cap \left( \text{ext}\, S \right)^c && \text{\ex{DeMorgansLaws}{De Morgan's laws}} \\
                       &= \left( \text{int}\,S \right)^c \cap \left( \overline{S}^c \right)^c && \text{\ex{Exteriors}{Exteriors}} \\
            &= \left( \text{int}\,S \right)^c \cap \overline{S} \\
            &= \overline{S} \setminus \text{int}\,S
        \end{align*}
    \end{proof}

    \begin{characterisation*} \label{characterisation}
        $p \in \partial S \iff \forall A \in \mathcal{N}_p, S \cap A \ne \emptyset \land S^c \cap A \ne \emptyset$\footnote{\label{fn:neighbourhoods}\ex{Neighbourhoods}{Neighbourhoods}}.
    \end{characterisation*}
    \begin{proof}
        \phantom{}\vspace{-2em}
        \begin{align*}
            p \in \partial S &\iff p \in (\inte S \cup \ext S)^c \\
            p \notin \partial S &\iff p \in \inte S \cup \ext S \\
            &\iff p \in \inte S \lor p \in \ext S \\
            &\iff \left( \exists A \in \mathcal{N}_p : A \subseteq S \right) \lor \left( \exists A \in \mathcal{N}_p : S \cap A = \emptyset \right) && (1) \\
            &\iff \exists A \in \mathcal{N}_p : A \subseteq S \lor S \cap A = \emptyset \\
            &\iff \exists A \in \mathcal{N}_p : S \cap A = \emptyset \lor S^c \cap A = \emptyset \\
            p \in \partial S &\iff \forall A \in \mathcal{N}_p, S \cap A \ne \emptyset \land S^c \cap A \ne \emptyset,
        \end{align*}
        where (1) is by the \ex[characterisation]{Interiors}{characterisation of interiors} and \ex[characterisation]{Exteriors}{characterisation of exteriors}. $\qed$
    \end{proof}

    \begin{theorem} \label{complement}
        $\partial (S^c) = \partial S$
    \end{theorem}
    \begin{proof}
        \phantom{}\vspace{-2em}
        \begin{align*}
            \partial (S^c) &= \left( \text{int}\,S^c \cup \text{ext}\, S^c \right)^c \\
                           &= \left( \text{ext}\,S \cup \text{int}\, S \right)^c && \text{\ex[complement]{Exteriors}{exteriors, theorem} \exref[complement]{Exteriors}} \\
            &= \partial S. \qed
        \end{align*}
    \end{proof}

    \begin{theorem} \label{open}
        $S \in \mathcal{T} \iff S \cap \partial S = \emptyset$ ($\partial S \subseteq S^c$).
    \end{theorem}
    \begin{proof}
        ($\implies$)
        Suppose $S \in \mathcal{T}$. By \textit{definition},
        \begin{align*}
            \partial S &= \overline{S} \setminus \inte S \\
            &= \overline{S} \setminus S. && \text{\ex[open]{Interiors}{interiors, theorem} \exref[open]{Interiors}}
        \end{align*}
        Thus, $S \cap \partial S = \emptyset$.

        ($\impliedby$)
        Let $p \in S$. Suppose that for all $A \in \mathcal{N}_p$\footref{fn:neighbourhoods}, $A \not\subseteq S$.
        Then there muxt exist $z \in A$ such that $z \in S^c$.
        Thus, by the \hyperref[characterisation]{characterisation of boundaries}, $p \in \partial S$.
        Thus, $p \in S \cap \partial S$, so $S \cap \partial S \ne \emptyset$.

        \ex{ConversesInversesAndContrapositives}{Contraposing} this, supposing $S \cap \partial S = \emptyset$ gives that $\forall p \in S, \exists A \in \mathcal{N}_p : A \subseteq S$\footref{fn:neighbourhoods}.

        Thus, by the \ex[characterisation]{Interiors}{characterisation of interiors}, $S \subseteq \inte S$.
        Together with the identity $\inte S \subseteq S$, this gives $S = \inte S$, and by \ex[open]{Interiors}{interiors, theorem}  \exref[open]{Interiors}, $S \in \mathcal{T}$. $\qed$
    \end{proof}

    \begin{theorem} \label{closed}
        $S^c \in \mathcal{T} \iff \partial S \subseteq S$.
    \end{theorem}
    \begin{proof}
        Substituting $S$ for $S^c$ in \hyperref[open]{theorem} \ref{open} gives
        $$ S^c \in \mathcal{T} \iff \partial (S^c) \subseteq (S^c)^c \iff \partial S \subseteq S, $$
        as $\partial S^c = \partial S$ by \hyperref[complement]{theorem} \ref{complement}. $\qed$
    \end{proof}
\end{document}
